\documentclass[12pt, a4paper, parskip=half]{scrartcl}

% ==========================================
% QUALITY OF LIFE PAKETE & EINSTELLUNGEN
% ==========================================
\usepackage[utf8]{inputenc}
\usepackage[T1]{fontenc}
\usepackage[ngerman]{babel}
\usepackage{lmodern}
\usepackage{amsmath, amssymb}
\usepackage{graphicx}
\usepackage{geometry}
\geometry{a4paper, left=2.5cm, right=2.5cm, top=2.5cm, bottom=2.5cm}
\usepackage{setspace}
\usepackage{csquotes}
\usepackage{booktabs}
\usepackage[hidelinks]{hyperref}

% ==========================================
% METADATEN FÜR DIE TITELSEITE
% ==========================================
\title{Kalibrierung von 2D-Kameras}
\subtitle{Seminararbeit im Modul Sensoren \& Aktoren für Robotersysteme}

% Formatierung der Autoren in einer Tabelle für saubere Ausrichtung
\author{
  \begin{tabular}{c p{1cm} c}
    \textbf{Nico Sander} & & \textbf{Marlon Leitschuh} \\
    \small \href{mailto:sani1019@h-ka.de}{sani1019@h-ka.de} & & \small \href{mailto:lema1046@h-ka.de}{lema1046@h-ka.de} \\
    \small Matrikelnr.: 74808 & & \small Matrikelnr.: 103626
  \end{tabular}
}

\date{\today}

\begin{document}

% Titelseite
\maketitle
\vfill
\begin{center}
    \small
    Betreuer: Prof. Dr. Daniel Braun \\
    Abgabedatum: [Datum]
\end{center}
\thispagestyle{empty}
\newpage

% Inhaltsverzeichnis
\tableofcontents
\newpage

\onehalfspacing

% ==========================================
% INHALT (Struktur aus Struktur_Seminararbeit.pdf)
% ==========================================

\section{Einleitung: Das 2D-Kamerasystem als Messsensor}

In der heutigen Industrielandschaft, die maßgeblich durch die vierte industrielle Revolution und die zunehmende Autonomie von Robotersystemen geprägt ist, hat sich die Rolle der 2D-Kamera grundlegend gewandelt. Sie fungiert nicht mehr bloß als Instrument zur visuellen Dokumentation oder zur einfachen Objekterkennung, sondern hat sich zu einem hochpräzisen optischen Messsensor entwickelt.[1, 2] In einem Masterstudium der Robotik und Künstlichen Intelligenz ist das Verständnis der Kamerasensorik von zentraler Bedeutung, da die visuelle Wahrnehmung oft die einzige Schnittstelle zwischen der digitalen Steuerungslogik und der physischen Realität darstellt.

\subsection{Grundlagen der optischen Projektion}

Eine Kamera ist im Kern ein optisches System, das die komplexe Geometrie einer dreidimensionalen Szene auf eine zweidimensionale Bildebene projiziert.[1] Dieser Prozess der Dimensionsreduktion ist mit einem Informationsverlust verbunden, da die Tiefeninformation primär verloren geht. Um diesen Prozess umkehrbar zu machen und metrische Daten aus den aufgenommenen Bildern zu extrahieren, ist eine mathematische Beschreibung dieser Abbildung unerlässlich.[3, 4] Die physikalische Grundlage bildet dabei das Licht, welches von Objekten reflektiert wird, durch ein Linsensystem tritt und schließlich auf einem Halbleitersensor (CMOS oder CCD) registriert wird.

Die Modellierung dieser Lichtwege folgt in der Regel dem Prinzip der Zentralprojektion, wobei alle Lichtstrahlen durch ein einzelnes Projektionszentrum verlaufen.[2, 3] In der Realität weichen industrielle Kamerasysteme jedoch von diesem idealisierten Modell ab. Fertigungstoleranzen bei der Montage der Sensoren, Unvollkommenheiten in der Linsenkrümmung und die präzise Ausrichtung der optischen Achse führen dazu, dass jedes Kamerasystem individuelle Charakteristika aufweist.[3, 5]

\subsection{Definition der Kamerakalibrierung}

Die Kamerakalibrierung wird definiert als der methodische Prozess zur Bestimmung der mathematischen Parameter, welche die Abbildung der 3D-Welt auf die 2D-Sensorebene beschreiben.[1, 6] Man unterscheidet hierbei zwischen der Bestimmung der inneren Geometrie des Sensors und der Optik (intrinsische Parameter) sowie der Bestimmung der Position und Orientierung der Kamera relativ zu ihrer Umgebung (extrinsische Parameter).[1, 7] Erst durch die Kenntnis dieser Parameter können Bildpunkte (Pixel) in reale Koordinaten (Millimeter oder Meter) umgerechnet werden, was die Basis für jede messende Anwendung darstellt.[1, 8]

\subsection{Zielsetzung des Kalibrierungsprozesses}

Das primäre Ziel der 2D-Kalibrierung ist die Minimierung des sogenannten Reprojektionsfehlers. Dies bedeutet, dass ein mathematisch berechneter Punkt in der Bildebene so exakt wie möglich mit dem tatsächlichen physikalischen Abbild auf dem Sensor übereinstimmen muss.[3, 9] Dabei verfolgt die Kalibrierung drei wesentliche Unterziele:
\begin{enumerate}
    \item Die Ermittlung der intrinsischen Kamerawerte wie Brennweite und Hauptpunktverschiebung.[1, 7]
    \item Die rechnerische Korrektur von optischen Verzeichnungen, um ein \enquote{ideales} Bild ohne Linsenfehler zu erzeugen.
    \item Die Bestimmung der Lage der Kamera im Raum, um eine Interaktion zwischen Roboterarmen und visuell detektierten Objekten zu ermöglichen.
\end{enumerate}

\subsection{Relevanz in der Robotik und KI}

Die Relevanz einer präzisen Kalibrierung ist in der modernen Robotik omnipräsent. Kleine Fehler in den Kameraparametern führen bei der Transformation in den 3D-Raum zu Fehlern, die proportional zur Entfernung des Objekts wachsen.[1, 3] Beispielsweise kann eine Abweichung der Brennweite von nur einem Prozent bei einer Entfernung von einem Meter bereits zu einem Lokalisierungsfehler von mehreren Zentimetern führen – eine Toleranz, die für präzise Pick-and-Place-Aufgaben oder die automatisierte Montage völlig inakzeptabel ist.[1]

Darüber hinaus benötigen moderne Algorithmen der Künstlichen Intelligenz, insbesondere im Bereich des maschinellen Sehens (Computer Vision), kalibrierte Eingangsdaten, um robuste Vorhersagen treffen zu können. Ein neuronales Netz, das auf die Erkennung von Defekten trainiert wurde, kann durch Linsenverzerrungen am Bildrand getäuscht werden, wenn diese nicht vorab rechnerisch kompensiert wurden. Somit bildet die Kalibrierung das Fundament, auf dem alle weiteren Schritte der Bildverarbeitung und Robotersteuerung aufbauen.[1, 6]

% Platzhalter für Abbildung: Projektionsmodell
\begin{figure}[h!]
    \centering
    \fbox{Abbildung: Darstellung der Projektion einer 3D-Szene auf eine 2D-Bildebene}
    \caption{Schematische Darstellung der Zentralprojektion und der Koordinatensysteme.[1, 3]}
    \label{fig:projection_model}
\end{figure}

\section{Das Kameramodell: Parameterarten}

Um die Abbildungseigenschaften einer Kamera rechnerisch handhabbar zu machen, wird ein mathematisches Modell verwendet, das die physikalischen Vorgänge abstrahiert. Dieses Modell wird durch verschiedene Parameterarten definiert, die in der Fachliteratur meist in intrinsische und extrinsische Parameter unterteilt werden.[1, 7]

\subsection{Intrinsische Parameter (Innere Orientierung)}

Die intrinsischen Parameter beschreiben die interne Geometrie und die optischen Eigenschaften der Kamera-Objektiv-Einheit. Sie sind fest mit der Hardware verbunden und ändern sich – sofern kein Zoom-Objektiv verwendet wird – nach der Kalibrierung nicht mehr.[7, 10]

\subsubsection{Die effektive Brennweite ($f_x, f_y$)}

Die Brennweite stellt den physikalischen Abstand zwischen dem optischen Zentrum der Linse und dem Bildsensor dar. In der digitalen Bildverarbeitung wird dieser Wert jedoch nicht in Millimetern, sondern in Pixeleinheiten angegeben.[10] Da die einzelnen Pixel auf einem Sensor nicht immer perfekt quadratisch sind, wird die Brennweite oft getrennt für die horizontale ($f_x$) und vertikale ($f_y$) Achse berechnet.[1, 7] Dies berücksichtigt den Skalierungsfaktor der Sensorelemente und stellt sicher, dass geometrische Formen wie Kreise im Bild nicht als Ellipsen interpretiert werden.[7, 10]

\subsubsection{Der Bildhauptpunkt ($c_x, c_y$)}

Der Bildhauptpunkt, auch Principal Point genannt, ist der Punkt, an dem die optische Achse den Bildsensor senkrecht schneidet.[7, 10] Idealerweise sollte dieser Punkt exakt im geometrischen Zentrum des Pixelsensors liegen (z. B. bei 960/540 bei einer Full-HD-Auflösung). In der Praxis führen jedoch minimale Ungenauigkeiten bei der Montage des Sensors im Gehäuse dazu, dass der Hauptpunkt um einige Pixel verschoben ist.[1, 5, 7] Die Kenntnis dieses Punktes ist entscheidend, da alle radialen Verzerrungsmodelle auf diesem Zentrum basieren.[8, 9]

\subsubsection{Skalierungsfaktoren und Scherung ($s$)}

Zusätzlich zu den Brennweiten werden Skalierungsfaktoren für die Pixelgeometrie berücksichtigt. Der Scherungsparameter ($s$ oder Skew) beschreibt, ob die x- und y-Achsen des Sensors perfekt rechtwinklig zueinander stehen.[1, 7] Bei modernen CMOS-Sensoren ist dieser Wert aufgrund der hochpräzisen Fertigungsverfahren in der Regel vernachlässigbar klein oder Null.[7, 11] Dennoch bleibt er Teil der allgemeinen Kameramatrix, um auch bei extremen Weitwinkelobjektiven oder älteren Scannersystemen die mathematische Korrektheit zu wahren.[2, 10]

\subsubsection{Die Kameramatrix $K$}

Alle intrinsischen Parameter werden kompakt in einer $3 \times 3$-Matrix, der sogenannten Kameramatrix $K$, zusammengefasst [1, 7]:
\[
K = \begin{bmatrix} f_x & s & c_x \\ 0 & f_y & c_y \\ 0 & 0 & 1 \end{bmatrix}
\]
Diese Matrix ist das Herzstück der intrinsischen Kalibrierung und ermöglicht die Transformation von Koordinaten im Kamerakoordinatensystem in Pixelkoordinaten.[1, 7]

\subsection{Extrinsische Parameter (Äußere Orientierung)}

Während die intrinsischen Parameter das \enquote{Innenleben} der Kamera beschreiben, definieren die extrinsischen Parameter die Lage und Ausrichtung des Kamerasystems relativ zum Einsatzort oder einem Weltkoordinatensystem.[1, 7] Dies ist vergleichbar mit der Angabe von Standort und Blickrichtung einer Person im Raum.

\subsubsection{Translation ($T$)}

Der Translationsvektor beschreibt die Verschiebung des Kameraursprungs (optisches Zentrum) zum Ursprung des Weltkoordinatensystems. Er besteht aus drei Parametern ($t_x, t_y, t_z$), welche die Position in den Raumrichtungen definieren.[1, 9] In Roboteranwendungen entspricht dies oft dem Abstand zwischen der Roboterbasis und dem Montagepunkt der Kamera.

\subsubsection{Rotation ($R$)}

Die Rotation beschreibt die Ausrichtung der Kamera im Raum. Hierfür sind drei Freiheitsgrade erforderlich, die meist als Rotationswinkel um die drei Achsen (Roll, Pitch, Yaw) angegeben werden.[9] Mathematisch wird dies durch eine $3 \times 3$ Rotationsmatrix $R$ dargestellt, welche sicherstellt, dass die Vektoren orthogonal bleiben.[1, 5] Die extrinsischen Parameter werden oft in einer homogenen Transformationsmatrix zusammengefasst, um 3D-Weltpunkte effizient in das lokale Koordinatensystem der Kamera zu überführen.[1, 3]

\begin{table}[h!]
    \centering
    \begin{tabular}{lll}
        \toprule
        Parameterart & Komponenten & Bedeutung \\
        \midrule
        \textbf{Intrinsisch} & $f_x, f_y, c_x, c_y, s$ & Interne Optik, Brennweite, Bildmitte \\
        \textbf{Extrinsisch} & $R (3 \times 3), T (3 \times 1)$ & Position und Orientierung im Weltraum \\
        \textbf{Verzerrung} & $k_1, k_2, p_1, p_2$ & Korrekturkoeffizienten für Linsenfehler \\
        \bottomrule
    \end{tabular}
    \caption{Zusammenfassung der wesentlichen Kameraparameter.[1, 7, 9, 10]}
\end{table}

\section{Physikalische Herausforderung: Linsenverzerrungen}

Kein optisches System ist perfekt. Die Verwendung von Glaslinsen zur Lichtbrechung führt zwangsläufig zu geometrischen Abweichungen, da Lichtstrahlen je nach Einfallswinkel unterschiedlich stark gebrochen werden. Diese Abweichungen werden als Linsenverzerrungen bezeichnet und müssen im Rahmen der Kalibrierung korrigiert werden, um metrische Genauigkeit zu erreichen.[1, 3]

\subsection{Ursache der Verzerrungen}

Die Hauptursache für Verzerrungen liegt in der sphärischen Form der Linsenoberflächen und der nicht idealen Ausrichtung der optischen Komponenten während der Fertigung.[1, 5] Während ein theoretisches Lochkameramodell Lichtstrahlen geradlinig passieren lässt, erzeugen reale Linsen eine Krümmung der Abbildung, die besonders zu den Bildrändern hin zunimmt.[3, 8]

\subsection{Radiale Verzerrung}

Die radiale Verzerrung ist die dominanteste Form des geometrischen Fehlers. Sie tritt auf, wenn der Vergrößerungsfaktor des Objektivs vom radialen Abstand zur optischen Achse abhängt.[3, 8]

\subsubsection{Tonnenförmige Verzerrung (Barrel Distortion)}

Bei der tonnenförmigen Verzerrung nimmt die Vergrößerung zu den Bildrändern hin ab. Dies führt dazu, dass gerade Linien im Außenbereich des Bildes wie die Seiten eines Fasses nach außen gebogen erscheinen.[1, 3] Dieser Effekt tritt typischerweise bei Weitwinkelobjektiven auf, da hier ein extrem großes Sichtfeld auf einen vergleichsweise kleinen Sensor komprimiert werden muss. In der Robotik ist dies bei Navigationskameras häufig zu beobachten, die ein breites Umfeld erfassen müssen.

\subsubsection{Kissenförmige Verzerrung (Pincushion Distortion)}

Umgekehrt nimmt bei der kissenförmigen Verzerrung die Vergrößerung zu den Bildrändern hin zu. Linien biegen sich zur Bildmitte hin ein, ähnlich wie ein eingedrücktes Kissen. Dieser Effekt ist oft bei Teleobjektiven zu finden, wenn das Sichtfeld des Objektivs kleiner ist als die tatsächliche Fläche des Bildsensors.[1]

Mathematisch wird die radiale Verzerrung durch Korrekturkoeffizienten ($k_1, k_2, k_3$) beschrieben, welche die Pixelpositionen basierend auf ihrem Abstand zum Hauptpunkt verschieben.[8, 9, 10]

\subsection{Tangentiale Verzerrung}

Die tangentiale Verzerrung entsteht durch mechanische Unvollkommenheiten, wenn die einzelnen Linsenelemente innerhalb des Objektivs oder die Linse zum Sensor nicht exakt parallel ausgerichtet sind. Dies führt zu einer leichten \enquote{Schrägstellung} der Abbildung.[1] In der industriellen Praxis ist dieser Fehler bei hochwertigen Objektiven oft vernachlässigbar klein, wird jedoch bei hochpräzisen Messaufgaben durch die Parameter $p_1$ und $p_2$ mitberücksichtigt.[1, 7, 10]

\subsection{Lösung durch Kalibrierung}

Die Kalibrierung ermittelt diese Korrekturkoeffizienten, indem sie die Differenz zwischen der verzerrten Beobachtung und der idealen projektiven Geometrie minimiert.[1, 9] Das Ergebnis ist ein Korrekturalgorithmus, der jedes aufgenommene Bild \enquote{entzerren} kann. Ein so behandeltes Bild verhält sich mathematisch wie ein Bild einer idealen Lochkamera, was die anschließende Berechnung von Distanzen und Posen massiv vereinfacht.[1, 3]

% Platzhalter für Abbildung: Verzerrungseffekte
\begin{figure}[h!]
    \centering
    \fbox{Abbildung: Gegenüberstellung von tonnenförmiger und kissenförmiger Verzeichnung}
    \caption{Visualisierung der radialen Verzeichnungstypen im Vergleich zum idealen Gitter.}
    \label{fig:distortions}
\end{figure}

\section{Geometrische Kalibrierverfahren}

Um die beschriebenen intrinsischen und extrinsischen Parameter zu bestimmen, wurden verschiedene mathematische Verfahren entwickelt. Diese nutzen in der Regel bekannte geometrische Strukturen, um die Abbildungsfehler der Kamera zu berechnen.[1, 3]

\subsection{Verwendung von Kalibriermustern}

Die zuverlässigste Methode zur Kamerakalibrierung ist die Verwendung spezieller Kalibriermuster. Dabei handelt es sich meist um plane Objekte mit hochpräzisen, kontrastreichen Mustern wie Schachbrettern oder Kreisgittern. Die Kamera nimmt dieses Muster aus verschiedenen Positionen und Orientierungen auf. Da die physikalischen Abstände auf dem Muster (z. B. 30 mm pro Quadratseite) exakt bekannt sind, kann der Computer die Abweichungen im Bild analysieren.

\subsection{Methode nach Tsai}

Die 1987 von Roger Tsai vorgestellte Methode ist ein klassisches photogrammetrisches Verfahren. Es handelt sich um eine zweistufige Technik zur Minimierung des Fehlers in der Bildebene.[3, 9]

\begin{enumerate}
    \item \textbf{Stufe 1}: Schätzung der extrinsischen Parameter (Rotation und Teile der Translation) sowie des Skalierungsfaktors durch lineare Gleichungssysteme.[3, 9]
    \item \textbf{Stufe 2}: Exakte Berechnung der Brennweite, der z-Translation und der radialen Verzerrung mittels nicht-linearer Optimierung (z. B. Gradientenabstieg).[3, 9]
\end{enumerate}

Tsai benötigt mindestens 7 bekannte Weltkoordinatenpunkte, wobei eine höhere Anzahl die Genauigkeit signifikant steigert.[1, 9] Ein Nachteil der Tsai-Methode ist die Anforderung an ein sehr präzise gefertigtes 3D-Kalibrierobjekt oder eine sehr genaue Positionierung, was in der Praxis oft aufwendig ist.[12]

\subsection{Methode nach Zhengyou Zhang}

Das Verfahren nach Zhang (entwickelt 1998/2000) ist heute der Standard in der industriellen Bildverarbeitung und in Bibliotheken wie OpenCV implementiert.[1, 1, 13]

Der große Vorteil von Zhangs Methode ist ihre Flexibilität: Sie benötigt lediglich ein flaches Schachbrettmuster, das einfach ausgedruckt und auf eine ebene Fläche geklebt werden kann.[1, 12, 13] Im Gegensatz zu Tsai muss die Position des Musters im Weltkoordinatensystem nicht vorab bekannt sein.[1, 13] Die Software detektiert die Eckpunkte des Schachbretts in den Bildern und nutzt die Tatsache aus, dass die Punkte in der Realität auf einer Ebene liegen ($Z=0$), um die intrinsischen Parameter zu extrapolieren. Durch die Aufnahme von etwa 10 bis 20 verschiedenen Posen des Schachbretts wird eine sehr hohe Genauigkeit erreicht, die für die meisten Robotikanwendungen absolut ausreichend ist.[1, 10, 12]

\subsection{Selbstkalibrierung}

Die Selbstkalibrierung ist ein fortgeschrittenes Verfahren, bei dem keine speziellen Muster benötigt werden. Stattdessen werden Merkmale der natürlichen Umgebung (Kanten, Ecken von Möbeln oder Gebäuden) genutzt, während die Kamera bewegt wird. Durch die Analyse der Kamerabewegung in einer Videosequenz können die Parameter geschätzt werden.[1, 2]

\begin{table}[h!]
    \centering
    \begin{tabular}{llll}
        \toprule
        Merkmal & Tsai-Methode & Zhang-Methode & Selbstkalibrierung \\
        \midrule
        \textbf{Bedarf} & 3D-Kalibrierobjekt & 2D-Flachmuster & Natürliche Merkmale \\
        \textbf{Bilderanzahl} & Wenige (oft eins) & Viele (10-20) & Kontinuierlicher Stream \\
        \textbf{Genauigkeit} & Mittel & Hoch & Geringer \\
        \textbf{Komplexität} & Hoch (Hardware) & Gering (Software) & Sehr hoch (Algorithmus) \\
        \bottomrule
    \end{tabular}
    \caption{Vergleich der gängigen geometrischen Kalibrierverfahren.[1, 12, 13]}
\end{table}

\subsection{Praktische Umsetzung in ROS und OpenCV}

In der Robotik wird häufig das Robot Operating System (ROS) eingesetzt. ROS bietet ein standardisiertes Paket zur Kamerakalibrierung, das auf der Zhang-Methode basiert.[1] Der Benutzer führt das Programm aus, hält ein Schachbrettmuster in das Sichtfeld der Kamera und bewegt es so lange, bis genügend Daten für alle Bildbereiche gesammelt wurden. Nach erfolgreicher Berechnung speichert ROS die Parameter in einer YAML-Datei, die automatisch von anderen Knoten (z. B. für die Objekterkennung) eingelesen wird, um entzerrte Bilder zu generieren.[1]

\section{Radiometrische Kalibrierung}

Während die geometrische Kalibrierung sicherstellt, dass die Strukturen im Bild an der richtigen Stelle sitzen, befasst sich die radiometrische Kalibrierung mit der Korrektur von Helligkeiten, Farben und Sensorfehlern.

\subsection{Fokus der radiometrischen Korrektur}

Das Ziel ist es, ein Bild zu erhalten, bei dem die Helligkeitswerte der Pixel direkt proportional zur tatsächlich einfallenden Lichtmenge sind.[1, 14] Dies ist besonders wichtig für die automatisierte Qualitätskontrolle, bei der Farbunterschiede oder Oberflächenbeschaffenheiten beurteilt werden müssen.

\subsection{Defekte Pixel}

Kein Kamerasensor ist perfekt. Aufgrund von Fertigungsfehlern oder Alterung können einzelne Pixel fehlerhafte Werte liefern. Man unterscheidet:
\begin{itemize}
    \item \textbf{Dead Pixel}: Pixel, die immer schwarz bleiben (kein Signal).[1, 8]
    \item \textbf{Hot Pixel}: Pixel, die aufgrund von Dunkelstrom immer weiß oder hell leuchten, auch wenn kein Licht einfällt.[1, 8]
\end{itemize}

Diese Fehler werden durch die Aufnahme von \textbf{Dunkelbildern} (Objektiv abgedeckt) und \textbf{Weißbildern} (homogen beleuchtete Fläche) identifiziert.[1, 8] In der Software werden diese Pixel dann durch Mittelwertbildung ihrer Nachbarn ersetzt oder schlichtweg ignoriert.[1, 8, 15]

\subsection{Randlichtabfall (Vignettierung)}

Vignettierung bezeichnet den Effekt, dass ein Bild zu den Rändern hin dunkler wird, obwohl das Motiv gleichmäßig hell ist. Dies hat zwei Ursachen:
\begin{enumerate}
    \item \textbf{Optische Vignettierung}: Das Objektivgehäuse schattet schräg einfallende Lichtstrahlen ab.[15]
    \item \textbf{Sensor-Vignettierung}: Die Mikrolinsen auf dem CMOS-Sensor fangen Licht am Rand weniger effizient ein als im Zentrum.
\end{enumerate}

Die Korrektur erfolgt durch die \textbf{Shading-Korrektur} (auch Flat-Field-Korrektur genannt). Dabei wird das Bild mit einem invertierten Weißbild verrechnet, um die Helligkeit über das gesamte Sichtfeld zu homogenisieren.[1, 14] Moderne Industriekameras können diese Korrektur oft direkt in der Hardware (FPGA) durchführen.[14, 15]

\subsection{Farbkalibrierung und Weißabgleich}

Für Anwendungen, die Farbtreue erfordern (z. B. Klassifikation von Lebensmitteln oder Lackprüfungen), ist eine Farbkalibrierung notwendig. Hierbei wird eine standardisierte Farbtafel (z. B. Color Checker) fotografiert. Die Abweichung der aufgenommenen RGB-Werte zu den Referenzwerten der Tafel wird genutzt, um eine Korrekturmatrix zu berechnen.

Häufig reicht jedoch ein einfacher \textbf{Weißabgleich} aus. Dabei wird festgelegt, welcher Farbwert als neutrales Weiß interpretiert werden soll. Da Lichtquellen unterschiedliche Farbtemperaturen haben (warmes Glühlampenlicht vs. kaltes Tageslicht), verhindert der Weißabgleich einen Farbstich im Bild und stellt sicher, dass weiße Objekte auch im digitalen Bild weiß erscheinen.[1, 15, 16]

\begin{table}[h!]
    \centering
    \begin{tabular}{lll}
        \toprule
        Korrekturtyp & Verfahren & Ziel \\
        \midrule
        \textbf{Defekte Pixel} & Dunkelbild/Maskierung & Entfernung von Bildfehlern \\
        \textbf{Vignettierung} & Flat-Field-Korrektur & Homogene Helligkeitsverteilung \\
        \textbf{Farbabgleich} & Farbtafel/Korrekturmatrix & Farbtreue Wiedergabe \\
        \textbf{Weißabgleich} & Referenzwert-Anpassung & Neutralisierung der Lichtquelle \\
        \bottomrule
    \end{tabular}
    \caption{Verfahren der radiometrischen Kalibrierung.[1, 8, 14, 15]}
\end{table}

\section{Anwendung und industrielle Praxis}

Die erfolgreiche Implementierung eines Kamerasystems in der Industrie hängt nicht nur von der Kalibrierung, sondern auch von der passenden Hardware-Wahl und den Umgebungsbedingungen ab.

\subsection{Anforderungen an Kamera und Optik}

Bei der Auswahl muss zwischen \textbf{Rolling Shutter} (günstiger, aber Verzerrungen bei Bewegung) und \textbf{Global Shutter} (präziser für bewegte Roboter) abgewogen werden. Die Brennweite muss zum Arbeitsabstand passen; oft wird die Formel $f = (h \cdot D) / (h + H)$ verwendet, um die benötigte Brennweite basierend auf Objektgröße $H$ und Distanz $D$ zu berechnen.[1] Eine hochwertige Optik reduziert zudem den Kalibrieraufwand, da sie weniger Verzerrungen und chromatische Aberrationen (Farbfehler) erzeugt.[1]

\subsection{Anforderungen an die Umgebung}

Beleuchtung ist das kritischste Element. Ein homogener Hintergrund und diffuses Licht verhindern störende Reflexionen auf metallischen Bauteilen, welche die Kalibrierung oder Objekterkennung verfälschen könnten. In der automatisierten Qualitätskontrolle werden oft spezielle Lichtzelte oder Ringlichter eingesetzt, um konstante Bedingungen zu schaffen.[1, 14]

\section{Fazit}

Die 2D-Kamerakalibrierung ist kein bloßer theoretischer Exkurs, sondern das unverzichtbare Fundament für alle messenden und wahrnehmenden Systeme in der Robotik und KI. Während die geometrische Kalibrierung durch Verfahren wie die Zhang-Methode heute weitgehend automatisiert und benutzerfreundlich umsetzbar ist, bleibt die radiometrische Korrektur entscheidend für die Stabilität von Bildverarbeitungsalgorithmen unter realen industriellen Bedingungen.[1, 14]

Ein Master-Absolvent im Bereich Robotik muss verstehen, dass die Qualität der Daten das gesamte System limitiert: Selbst die fortschrittlichste KI oder der präziseste Roboterarm können keine Höchstleistungen erbringen, wenn das \enquote{Auge} des Systems – die Kamera – nicht korrekt kalibriert ist. Die Integration von Hardware-Wissen (Sensortypen, Linsenphysik) mit mathematischen Modellen (Kameramatrizen, Optimierungsverfahren) ist daher die Schlüsselkompetenz für eine erfolgreiche Prozessintegration in der modernen Produktion.[1, 6]

\end{document}